\documentclass[12pt]{article}

\usepackage[margin=1in]{geometry}
\usepackage{setspace}
\usepackage{natbib}
\setstretch{1}

\setlength{\parindent}{0pt}
\setlength{\parskip}{0.8em}

\title{This is a title}
\author{Wenqin Huang}
\date{\today}

\begin{document}

\maketitle


As low Earth orbit broadband networks expand rapidly, space launches have become increasingly indispensable to global connectivity, yet their growing environmental costs can no longer be treated as a secondary concern. Space activity is no longer occasional; it is becoming a rapidly expanding industrial sector: according to Ryan et al.'s 2022 research, its scale will leap from 350 billion in 2019 to over a trillion by 2040~\citep{ryan2022ozone,bushnell2018newspace}. This critical leap will probably cause a substantial intensification and amplification of its environmental impacts. Additionally, through rocket launching activities, multiple pollutants will be directly injected into the upper atmosphere, where pollution will be far more destructive, posing a direct threat to ozone~\citep{ryan2022ozone,kokkinakis2022pollution}. Moreover, this expansion is not only something that will happen in the future, but is already ongoing: the frequency of rocket launches increased substantially from 58 launches in 2003 to 329 launches in 2025~\citep{mcdowell2025space}. The pollution is not simply limited to Earth's ecosystem, but also influences Earth's orbits. Satellite constellations will also occupy Earth's orbits, create hundreds of thousands of pieces of space debris, and subsequently obstruct future satellite deployments~\citep{torres2020reusable}. Thus, the environmental impact of rocket launching could no longer be ignored, since the industry is rapidly expanding, and it influences the environment in multiple different ways.



However, LEO broadband constellations are not simply symbolic space activities, but important components of communication infrastructure. The rapid expansion of OneWeb and Starlink is motivated and triggered by the growth of global communication demand. The connectivity created by LEO broadband systems is also widely argued to be beneficial, especially in remote areas, where traditional internet service does not reach: LEO satellites provide them with internet connections, and therefore improve their social connections, and subsequently mental health \citep{osoro2025leo,sandau2010smallsat}. Based on the benefits they bring, space launches should not simply be halted. Thus, the true essential problem is "to what extent can technological strategies reduce these environmental costs and improve sustainability", rather than "whether people should continue launching or not".



To address the environmental costs associated with frequent space launches, various technological strategies have been proposed to improve the sustainability of launch systems. Among these, reusable launch vehicles, alternative propellants, and satellite miniaturization are often considered to be the most practical approaches. Their effects can be evaluated based on their ability to reduce emissions, minimize life-cycle environmental impacts, and remain effective under the rapid expansion of satellite deployment. However, the effectiveness and validity of these approaches are not essentially limited by technological capability, but are complicated by financial incentives, technological limitations, and the rebound effect that may offset or even reverse their environmental benefits. Therefore, although technological strategies are already sufficiently developed to reduce the environmental impact of space launches, their real-world effectiveness remains constrained by economic barriers, industrial inertia, and the rapid expansion of LEO megaconstellations, limiting their practical implementation.



Alternative propellants are the most direct technique to reduce environmental impact among the three approaches, since they change the fuel combustion products directly. Dallas et al. systematically reviewed the environmental impact brought by traditional propellants, and argue that solid fuel, which typically involves chemicals like ammonium perchlorate, may potentially damage the ozone the most; kerosene-based propellants mainly emit carbon dioxide (CO$_2$) and black carbon, which may intensify the greenhouse effect, but are normally less destructive than solid propellants~\citep{dallas2020review,brown2024inventory}. By contrast, alternative fuels are suggested to be more environmentally friendly. Liquid hydrogen is widely recommended as a candidate for future propulsion, since the emissions it generates contain mainly water vapor and have negligible ecological toxicity~\citep{dallas2020review,barato2023fuels}. To be more specific, Calabuig et al.'s study suggests that the carbon footprint of liquid hydrogen (alternative propellant) fleets is typically 2--8 times lower than that of methane-based (transitional propellant, cleaner than kerosene, but not as clean as hydrogen) vehicle fleets. Additionally, they further argue that, in most cases, choosing sustainable propellants will affect more than the vehicle's reusability: propellant choice is the main differentiation factor when evaluated under environmental criteria~\citep{calabuig2024reusable}. Recent research by Hai et al. suggests that environmentally friendly propellants not only reduce toxic products, but also lower the risks of transportation and operation~\citep{hai2026propellants}. These findings collectively suggest that alternative propellants provide a highly effective approach for reducing the environmental impact of rocket launches in principle. However, their real-world implementation remains obstructed by practical and economic challenges. The storage and transportation of alternative propellants are typically complex and expensive. Dallas et al. argue that liquid hydrogen is inherently unsuitable for storage: its density is lower, which means it requires a storage system that is significantly larger and more complex than traditional fuels; additionally, the boiling point of liquid hydrogen is considerably lower, which means the propellant will be more likely to vaporize---existing propellant storage tanks may not be suitable for storage~\citep{dallas2020review}. This indicates that the cleaner propellant may not be easier to store or use. Additionally, Calabuig et al. also argue that although the technology is proven beneficial, the cost is still significantly higher than that of traditional propellants, and the industry, including production, transportation, and storage, is also immature. Since more and more rocket launching activities are nowadays carried out by the private sector, rather than government-supported institutions, the higher cost may discourage companies from utilizing alternative propellants~\citep{calabuig2024reusable}. Therefore, the established aerospace industry was structured around traditional propellants like kerosene, and the transformation may not be financially favorable. In conclusion, alternative propellants offer a concrete environmentally friendly choice and emission reduction strategy, but their effect depends on the pace of industrial transformation and financial incentives, rather than technological limitations.



However, propellant choice alone does not fully determine the environmental impact of space launches. Beyond fuel composition, the design and operational characteristics of launch systems---particularly the use of reusable launch vehicles---also play a critical role. Reusable launch vehicles offer a promising solution to reduce the environmental costs of space launches by reducing manufacturing emissions and waste. Therefore, applying reusable rocket technology can reduce the manufacturing cost of launching, but it will not necessarily reduce the net environmental impact, Calabuig et al. argue~\citep{calabuig2024reusable}. On one hand, reusable vehicles can improve the material efficiency of fixed equipment and the fleet, and therefore reduce the launching cost financially. As Calabuig writes, the original idea of reusable rockets is to reduce material waste, and when evaluated under the criterion of launching impact per unit of mass, reusable vehicles can indeed improve both environmental and financial efficiency. However, Calabuig et al. also provide solid numerical evidence proving that the improvement by reusable rockets is limited: for liquid hydrogen fleets, the climate impact could not be reduced by more than 40\%; for methane fleets, the impact is less significant, around 20\%~\citep{calabuig2024reusable}. This trend indicates that the climate impact reduction is considerable, but not as transformative as alternative propellants. However, long-term assessments, like life-cycle evaluations, are also critical in evaluating the environmental impact of an industry~\citep{wilson2022lifecycle}. Under long-term evaluation, because such vehicles are more financially favorable, with the same budgets, companies may launch substantially more frequently and combust a significantly greater amount of propellants, which, as mentioned in the last paragraph, could amplify the environmental impact and carbon emissions to a greater extent, causing a rebound effect. The rebound effect of reusable vehicles is well illustrated by SpaceX, a private launch provider that has pioneered the large-scale adoption of reusable rocket technology, particularly through its Falcon 9 program. According to FAA environmental reviews, the annual Falcon launch cadence considered for SpaceX has risen steadily, from 20 launches per year at LC-39A and 50 at SLC-40 in the 2020 assessment, to 36 launches per year at Vandenberg in 2023, 50 in 2024, and 100 in the 2025 Vandenberg EIS \citep{faa2025falconprogram}, while it reached 165 launches in 2025 alone~\citep{mcdowell2025space}. These statistics confirm the leap in launch frequency under the implementation of reusable vehicle fleets. To summarize, although reusable launch vehicles can improve efficiency by reducing manufacturing impacts and lowering costs, their overall environmental benefits remain limited and could be offset, or even reversed, by increased launch frequency driven by economic incentives.



Satellite miniaturization improves the functional efficiency of individual satellites, but it also lowers the barrier to large-scale deployment. As Sandau et al. suggested, satellite miniaturization could improve the efficiency per unit of mass by reducing the mass while the function remains unchanged, and enable the satellite broadband system to cover globally~\citep{sandau2010smallsat}. Satellite miniaturization also makes satellites cheaper, easier to design, and easier to deploy. For environmental impact, the miniaturization of satellites means that, with the same rocket, more satellites could be launched and deployed; therefore, the average environmental cost of each satellite could be reduced. However, the miniaturization of satellites could also increase the scale of broadband constellations, while the increase in scale may not only intensify the environmental burden, but also occupy the orbit and obstruct future launches~\citep{kumaran2024manufacturing}. While megaconstellations enabled by miniaturization significantly improve broadband access in remote and underserved regions, their emissions intensity per subscriber remains substantially higher than that of terrestrial 4G systems, indicating a non-negligible environmental cost~\citep{osoro2025leo,kukreja2025leo}. This indicates that the benefits of satellite miniaturization are highly context-dependent: although efficiency is improved at the level of individual satellites, the expansion of constellation scale shifts the problem from per-unit optimization to total system impact.



After analyzing three technological strategies including alternative propellants, reusable rockets and satellite miniaturization, we can find that all three technologies are theoretically capable of reducing the environmental impact of satellite megaconstellations, but their effectiveness in real-world contexts is still largely limited. The key to addressing this problem is not the effectiveness of technology, but the complicated economic and industrial environment. Firstly, the market's incentive mechanism wasn't designed to optimize environmental goals: the reduction of cost would not reduce total launches, but would trigger more frequent spaceflight missions and deployments. To be more specific, the development of reusable rockets escalates spaceflight frequency by reducing its cost, and satellite miniaturization expands the megaconstellation of broadband networks by lowering per-satellite manufacturing cost. This scope expansion driven by technological advancements and efficiency improvements offsets their net environmental benefits, and may even trigger reverse effects. Secondly, the existing aerospace industry exhibits significant path dependence. For a long time, the rocket industry, including launchpad infrastructure, propellant storage systems, and related supply chains, has been established around traditional fuel systems. This means that although alternative propellants may offer better environmental performance, they may still face limitations in implementation and transformation due to the high infrastructure-reconstruction costs and industrial inertia. This inertia makes it hard to practically translate this technological effectiveness into large-scale adoption. Furthermore, to a great extent, the environmental impacts are derived from the expansion of scale, rather than solely caused by technological deficiency and inefficiency. In this case, it will be difficult to control the environmental burden by solely relying on technological advancements and optimizations. Therefore, this evidence and reasoning collectively suggest that the main factor limiting spaceflight sustainability is implementation and transformation, rather than technological limitations. This also means that, in order to make aerospace activities sustainable, institutional frameworks and policy tools, like limiting launch frequency and guiding industrial transition, are also required.



In conclusion, the rapid expansion of LEO broadband satellite constellations has intensified the environmental challenges associated with space launches, making sustainability a critical concern for the future of the space industry. While technological advancements, like alternative propellants, reusable rockets and satellite miniaturization, offer practical approaches to reduce environmental impacts, their effectiveness remains controversial and limited in practice. Alternative propellants provide the most direct way to reduce emissions, but face financial barriers and industrial inertia. Reusable launch vehicles improve per-launch efficiency and reduce costs financially, but their benefits are limited by the rebound effect driven by cost reductions. Similarly, satellite miniaturization enhances functional efficiency, but often leads to large-scale deployment that increases the total environmental burden.



These findings suggest that the primary limitation is not the lack of technological capability, but the implementation of these existing technologies. As long as economic incentives favor expansion and existing industrial systems do not transition, technological improvements alone are not likely to achieve significant reductions in environmental impact. Therefore, ensuring the sustainability of space activities requires not only technological innovation, but also structural changes in political and regulatory frameworks and economic incentives. Finally, sustainable space development depends on promoting the application of technological progress with system-level regulation, rather than relying on methodological advancements alone.

%\subsection{Why Mature Technologies Remain Underutilized}

%\subsection{Barriers to Adoption}

%\subsection{The Case for Regulatory Intervention}

%\section{Conclusion}

\newpage
\bibliographystyle{apalike}
\bibliography{references}


\end{document}
